\documentclass[11pt,a4paper]{article}

% --- Paquetes fundamentales ---
\usepackage[utf8]{inputenc}
\usepackage[spanish,es-tabla,es-noquoting]{babel}
\usepackage[margin=2.5cm]{geometry}
\usepackage{amsmath,amsfonts,amssymb}
\usepackage{booktabs}
\usepackage{tabularx}
\usepackage{array}
\usepackage{xcolor}
\usepackage{graphicx}
\usepackage{hyperref}
\usepackage{microtype}
\usepackage{enumitem}
\usepackage{tcolorbox}
\usepackage{caption}

% --- Configuración de enlaces y estilos ---
\hypersetup{
    colorlinks=true,
    linkcolor=blue!70!black,
    citecolor=blue!70!black,
    urlcolor=teal!70!black,
    pdfauthor={Equipo de Ingeniería Chiripa},
    pdftitle={Reporte Técnico - Spike 2: Matching Semántico}
}

\tcolorboxenvironment{quote}{
    colback=gray!5!white,
    colframe=gray!40!black,
    boxrule=0.5pt,
    left=8pt,
    right=8pt,
    top=6pt,
    bottom=6pt
}

\title{\textbf{Reporte Técnico --- Spike 2:}\\[0.3em]
\Large Evaluación, Calibración y Optimización del Motor de Matching Semántico}
\author{\textbf{Equipo de Ingeniería Chiripa} \\
Plataforma de Licitaciones Inteligentes (\textit{ProyectosYA})}
\date{Octubre 2026}

\begin{document}

\maketitle

\begin{abstract}
Este informe documenta exhaustivamente la investigación, experimentación y mejoras arquitectónicas realizadas en el \textbf{Spike 2} sobre el motor de recomendación y \textit{matching} semántico de Chiripa. A partir de una auditoría metodológica rigurosa sobre datos reales de Compra Ágil (Mercado Público), se eliminaron sesgos de fuga temporal y catálogo cerrado presentes en versiones preliminares. Evaluamos el comportamiento del \textit{baseline} puramente denso (\textit{BGE-M3}), el impacto de la incorporación de canales léxicos dispersos (\textit{BM25}), el aporte de señales estructuradas (región canónica y taxonomía UNSPSC), y la recalibración probabilística ordinal de la escala de compatibilidad mediante regresión logística multivariante ($R_c + B + C$). Los resultados demuestran que un enfoque híbrido triplica la calidad del ranking ($NDCG@10$ de $0{,}0517$ a $0{,}1579$, $p=0{,}002$), mientras que la calibración elevó el AUC de $0{,}58$ a $0{,}75$, permitiendo una detección certera de licitaciones de alta afinidad (precisión del $97{,}6\%$ en el umbral $\ge 70\%$).
\end{abstract}

\tableofcontents
\newpage

% =========================================================================
\section{Objetivo del Spike}
% =========================================================================

El propósito primordial del Spike 2 fue diagnosticar, medir objetivamente y optimizar el rendimiento del motor de búsqueda y recomendación de licitaciones para micro, pequeñas y medianas empresas (MiPymes). En Chiripa, una recomendación errónea o mal priorizada no solo reduce la interacción del usuario, sino que genera un costo de oportunidad severo: las MiPymes cuentan con tiempo limitado para revisar bases y presentar ofertas contractuales.

Los objetivos específicos del spike se desglosan en:

\begin{enumerate}[leftmargin=*]
    \item \textbf{Establecer un arnés de evaluación empírico sin fuga de datos (\textit{Leak-Free Benchmark}):} Diseñar un protocolo experimental con división temporal estricta y catálogo abierto con distractores reales de Mercado Público, evitando sobreajustes o métricas artificialmente infladas.
    \item \textbf{Evaluar cuantitativamente el pipeline denso inicial:} Medir con métricas estándar de Recuperación de Información (\textit{Information Retrieval, IR}) la capacidad del modelo vectorial \textit{BGE-M3} para posicionar compras adjudicadas a las empresas.
    \item \textbf{Explorar arquitecturas de recuperación híbridas y multietapa:} Analizar si la combinación de vectores densos con índices léxicos dispersos (\textit{sparse BM25}), señales estructuradas (región y categoría) y modelos de reordenamiento (\textit{cross-encoders rerankers}) produce mejoras con significancia estadística.
    \item \textbf{Calibrar la escala del porcentaje de compatibilidad ($[0\%, 100\%]$):} Corregir la compresión sistemática de puntajes que impedía a las licitaciones pertinentes alcanzar el umbral verde del sistema ($\ge 70\%$).
    \item \textbf{Definir la arquitectura productiva:} Proveer bases técnicas para el despliegue del canal léxico en Qdrant, la evaluación en sombra (\textit{Shadow Matching}) y la telemetría en línea.
\end{enumerate}

% =========================================================================
\section{Configuración Inicial del Pipeline de Matching}
% =========================================================================

Para que cualquier ingeniero de software comprenda el funcionamiento del sistema sin necesidad de especialización en procesamiento de lenguaje natural (NLP) o búsqueda semántica, en esta sección se formaliza cómo operaba el pipeline original.

\subsection{Conceptos Fundamentales}
\begin{itemize}[leftmargin=*]
    \item \textbf{Vector / Embedding:} Un vector es una lista ordenada de 1024 números decimales generada por una red neuronal (\textit{BGE-M3}). Este vector posiciona el texto en un espacio geométrico de alta dimensión: textos conceptualmente afines (por ejemplo, ``aseo industrial'' y ``limpieza hospitalaria'') quedan ubicados cerca espacialmente, incluso si no comparten palabras exactas.
    \item \textbf{Similitud Coseno:} Mide el ángulo entre dos vectores $\mathbf{u}$ y $\mathbf{v}$ de norma unitaria:
    \begin{equation}
        \text{sim}(\mathbf{u}, \mathbf{v}) = \frac{\mathbf{u} \cdot \mathbf{v}}{\|\mathbf{u}\| \|\mathbf{v}\|} \in [-1, 1]
    \end{equation}
    En la práctica, valores cercanos a 1 indican gran afinidad temática.
    \item \textbf{Base de Datos Vectorial (Qdrant):} Motor especializado que indexa millones de vectores para encontrar los $K$ vectores más cercanos a una consulta en milisegundos mediante índices HNSW (\textit{Hierarchical Navigable Small World}).
\end{itemize}

\subsection{El Algoritmo Inicial (P0) y la Generación del Ranking}
El flujo previo de recomendación operaba en tres etapas:

\begin{enumerate}[leftmargin=*]
    \item \textbf{Recuperación Bi-canal:}
    \begin{itemize}
        \item \textit{Canal 1 (Perfil completo):} Se calculaba el vector denso del perfil de la empresa (descripción + sectores) y se consultaba la colección \texttt{tenders} en Qdrant, recuperando las 50 licitaciones más afines.
        \item \textit{Canal 2 (Keywords vs. Partidas):} Se proyectaban las palabras clave individuales del proveedor contra las partidas específicas licitadas (\texttt{tender\_items}) mediante similitud MaxSim.
    \end{itemize}
    \item \textbf{Unión de Candidatas:} Se formaba un conjunto de hasta 50--80 licitaciones candidatas combinando ambos canales.
    \item \textbf{Puntuación Inicial por Combinación Lineal (\textit{FieldWeightingService}):}
    A cada licitación candidata se le calculaba una afinidad según la fórmula ad-hoc:
    \begin{equation}
        P_0 = 0{,}50 \cdot \text{Reranker} + 0{,}25 \cdot \text{Rubro} + 0{,}25 \cdot \text{Keywords} + \text{BonoLéxico}
    \end{equation}
    donde el \textit{Reranker} era un modelo cross-encoder (\textit{BGE-Reranker-v2-m3}) que evaluaba el par completo (proveedor, licitación).
\end{enumerate}

\subsection{Diagnóstico del Fracaso de la Configuración Inicial}
A pesar de su aparente sofisticación, este pipeline presentaba tres fallas críticas de diseño:

\begin{enumerate}[leftmargin=*]
    \item \textbf{Colapso del Bono Léxico:} El bono requería coincidencias de subcadenas literales idénticas. Diferencias triviales de plurales, acentuación o redacción hacían que casi siempre valiera 0.
    \item \textbf{Compresión de Escala:} Dado que los componentes de rubro y keywords rara vez daban 1 simultáneamente, el puntaje práctico máximo alcanzable rondaba el $30\%$--$40\%$. En pruebas de usuario real (por ejemplo, con la empresa de transportes \textit{Buses del Sur SpA}), servicios de traslado escolar directos recibían apenas un $25\%$ de compatibilidad, quedando excluidos del filtro verde ($\ge 70\%$).
    \item \textbf{Inestabilidad del Reranker INT8 por Cuantización Dinámica:} Al procesar lotes (\textit{batches}) de 50 elementos con modelos cuantizados en INT8, la cuantización recalibraba su rango según el contenido del lote entero. Así, el puntaje de la licitación $X$ cambiaba arbitrariamente según qué otras 49 licitaciones la acompañaban en la consulta.
\end{enumerate}

% =========================================================================
\section{Métricas de Evaluación: Definición, Cálculo e Interpretación}
% =========================================================================

Para aislar la subjetividad y medir con rigor científico, se adoptó el estándar de evaluación de \textit{Information Retrieval} y aprendizaje estadístico.

\subsection{Definición de Relevancia (Ground Truth)}
El conjunto de evaluación se construyó a partir de compras públicas reales de Mercado Público:
\begin{itemize}[leftmargin=*]
    \item \textbf{Relevancia Dura ($rel \in \{0, 2\}$):} $rel = 2$ si la orden de compra fue efectivamente adjudicada y ganada por el proveedor en una fecha \textbf{estrictamente posterior} a las órdenes que componen su perfil histórico. $rel = 0$ para cualquier otra licitación. Cada proveedor cuenta con exactamente 2 positivos futuros ($rel=2$) entre $1.176$ licitaciones del catálogo.
    \item \textbf{Relevancia Graduada con Juez Ciego ($rel \in \{0, 1, 2\}$):} Evaluada sobre un pool de $1.146$ pares evaluados por un LLM sin conocer si la orden fue adjudicada: $2$ (mismo producto exacto), $1$ (misma familia/rubro cotizable), $0$ (rubro no relacionado).
\end{itemize}

\subsection{Métricas de Ranking}

\subsubsection{NDCG@10 (Normalized Discounted Cumulative Gain a corte 10)}
Mide la calidad del orden posicional, premiando con mayor peso que los elementos relevantes aparezcan en los primeros lugares.
\begin{equation}
    DCG@K = \sum_{i=1}^{K} \frac{2^{rel_i} - 1}{\log_2(i + 1)}, \quad NDCG@K = \frac{DCG@K}{IDCG@K}
\end{equation}
donde $IDCG@K$ es el $DCG$ de la lista perfectamente ordenada. $NDCG \in [0, 1]$.
\begin{quote}
\textit{Interpretación:} Es la métrica reina del benchmark. Refleja qué tan optimizada está la primera pantalla visible por el usuario. Un incremento en NDCG@10 indica que las licitaciones pertinentes subieron a los primeros puestos.
\end{quote}

\subsubsection{MRR (Mean Reciprocal Rank)}
Evalúa la rapidez con la que el usuario encuentra su primer acierto relevante:
\begin{equation}
    MRR = \frac{1}{|Q|} \sum_{q=1}^{|Q|} \frac{1}{\text{rank}_{q}^{*}}
\end{equation}
donde $\text{rank}_{q}^{*}$ es la posición (1-indexed) de la primera licitación relevante para la consulta $q$. Si el primer acierto está en la posición 1, vale $1{,}0$; si está en la posición 2, vale $0{,}50$; si no aparece en el ranking evaluado, vale 0.

\subsubsection{Success@1 (Hit Rate en la Posición 1)}
Fracción binaria de consultas donde la recomendación que encabeza la lista (posición 1) es efectivamente relevante:
\begin{equation}
    \text{Success@1} = \frac{1}{|Q|} \sum_{q=1}^{|Q|} \mathbb{I}(rel_1 > 0)
\end{equation}

\subsubsection{P@10 (Precisión a los 10 primeros)}
Proporción de elementos relevantes dentro de los primeros diez resultados:
\begin{equation}
    P@10 = \frac{1}{10} \sum_{i=1}^{10} \mathbb{I}(rel_i > 0)
\end{equation}
\begin{quote}
\textbf{Nota técnica fundamental de interpretación:} En este benchmark, cada empresa tiene exactamente 2 licitaciones adjudicadas posteriores en el catálogo de $1.176$. Por lo tanto, el \textbf{techo teórico absoluto} de P@10 es $\frac{2}{10} = 0{,}20$ ($20\%$). Un $P@10 = 0{,}048$ representa un $24\%$ del techo máximo posible.
\end{quote}

\subsubsection{Recall@K \texorpdfstring{($K \in \{10, 50, 100, 200\}$)}{(K en 10, 50, 100, 200)}}
Fracción del total de licitaciones relevantes recuperadas dentro de los primeros $K$ candidatos:
\begin{equation}
    Recall@K = \frac{|\text{Relevantes en el top-}K|}{|\text{Total de relevantes para el usuario}|}
\end{equation}
\begin{quote}
\textbf{Techo de Rerankers:} El $Recall@100$ del primer estadio de búsqueda establece el límite insuperable de cualquier modelo de reordenamiento posterior. Si un acierto relevante no queda dentro de los 100 candidatos recuperados inicialmente, un reranker jamás podrá rescatarlo.
\end{quote}

\subsection{Inferencia y Significancia Estadística}
Para evitar validar mejoras espurias causadas por varianza aleatoria:
\begin{itemize}[leftmargin=*]
    \item \textbf{Prueba de Rangos Signados de Wilcoxon Pareado:} Evalúa si la diferencia de desempeño consulta por consulta entre dos estrategias difiere de cero sin asumir normalidad en los errores.
    \item \textbf{Corrección de Holm-Bonferroni:} Ajusta los valores $p$ para controlar la tasa de error por familia ante múltiples comparaciones.
    \item \textbf{Bootstrap Pareado (5.000 repeticiones):} Estima el intervalo de confianza al $95\%$ ($\text{IC}_{95\%}$) para el incremento medio $\Delta NDCG@10$. Una mejora es concluyente si su límite inferior es estrictamente positivo ($\text{IC}_{\text{inf}} > 0$).
\end{itemize}

% =========================================================================
\section{Resultados de Métricas sobre la Configuración Inicial}
% =========================================================================

La evaluación empírica del baseline puramente denso (\textbf{A: Denso BGE-M3}) sobre las 50 empresas de prueba arrojó resultados concluyentes sobre su insuficiencia operativa:

\begin{table}[htbp]
\centering
\small
\caption{Desempeño del Baseline Denso Inicial ($N=50$ consultas, Catálogo = $1.176$).}
\label{tab:baseline_inicial}
\begin{tabularx}{\textwidth}{lcccccccc}
\toprule
\textbf{Estrategia} & \textbf{NDCG@10} & \textbf{MRR} & \textbf{Succ@1} & \textbf{P@10} & \textbf{R@10} & \textbf{R@50} & \textbf{R@100} & \textbf{R@200} \\
\midrule
\textbf{A: Denso (BGE-M3)} & 0{,}0517 & 0{,}0675 & 0{,}0200 & 0{,}0160 & 0{,}0800 & 0{,}2400 & 0{,}3800 & 0{,}5300 \\
\bottomrule
\end{tabularx}
\end{table}

\subsection{Diagnóstico del Baseline Denso}
\begin{enumerate}[leftmargin=*]
    \item \textbf{Ceguera Temática Fina:} De las 50 consultas evaluadas, \textbf{solo 7 consultas lograron posicionar al menos un acierto en su top-10}. En el $86\%$ de los casos restantes, las licitaciones adjudicadas no aparecieron en la primera página.
    \item \textbf{Cuello de Botella Severo en Recuperación ($Recall@100 = 38\%$):} El $62\%$ de las compras públicas efectivamente ganadas por las empresas quedaron fuera de los primeros 100 resultados.
    \item \textbf{Causa Raíz:} Los textos de Mercado Público y Compra Ágil son breves, administrativos (frecuentemente titulados simplemente ``Orden de Compra generada'') y contienen especificaciones numéricas o códigos de producto. Un vector denso global resume el ``giro de negocio'' abstracto, pero fracasa al discriminar códigos de insumos específicos, marcas o servicios locales puntuales.
\end{enumerate}

% =========================================================================
\section{Espacio de Decisiones: Consideradas, Descartadas e Implementadas}
% =========================================================================

Durante el desarrollo del spike se investigaron múltiples hipótesis y componentes arquitectónicos. A continuación se detallan las alternativas exploradas y el fundamento técnico de su adopción o descarte.

\subsection{Decisiones Descartadas}

\begin{itemize}[leftmargin=*]
    \item \textbf{Descarte 1: Reranker BGE como solución al cuello de botella primario (Estrategia D).}
    \textit{Fundamento:} Aplicar el cross-encoder \textit{BGE-Reranker-v2-m3} sobre el top-100 no mejoró de forma estadísticamente significativa al híbrido estructural ($NDCG@10$ de $0{,}1175$ vs. $0{,}1579$, $p=0{,}063$, no significativo en métricas duras). Además, frente a un juez ciego, el reranker empeoró el ranking en $\Delta = -0{,}060$ ($p=0{,}005$). Dado que el reranker agrega un costo computacional no despreciable ($\sim 7$ segundos por consulta), se concluyó que el problema prioritario radicaba en el \textit{Recall} del primer estadio y no en el reordenamiento.
    
    \item \textbf{Descarte 2: Multi-consulta densa por producto histórico (Estrategia E).}
    \textit{Fundamento:} Generar una consulta densa individual por cada producto histórico y fusionar mediante el puntaje máximo aportó una ganancia marginal frente al híbrido léxico simple ($\Delta = +0{,}022$, $p=0{,}125$, no significativo) a cambio de multiplicar las llamadas de embeddings y búsquedas vectoriales por cada ítem del historial.
    
    \item \textbf{Descarte 3: Representaciones Léxicas de BGE-M3 (\textit{Sparse Lexical Weights}).}
    \textit{Fundamento:} BGE-M3 ofrece pesos dispersos nativos, pero su generación requiere la librería pesada \texttt{FlagEmbedding} con PyTorch y dependencias C++ no disponibles en las APIs comerciales ligeras ni deseables en los contenedores de producción de bajo consumo. Se optó por BM25 clásico con tokenización en la capa de aplicación.
    
    \item \textbf{Descarte 4: Inferencia de taxonomía UNSPSC para empresas sin historial (\textit{Cold Start} - Fase 2b).}
    \textit{Fundamento:} En la Fase 2 se evaluó si los proveedores recién registrados (sin órdenes previas) sufrían una brecha grave de calidad. La brecha en el pipeline híbrido fue estadísticamente insignificante ($\Delta NDCG@10 = -0{,}0048$, $p=0{,}483$, no significativo). Al intentar inferir categorías UNSPSC proyectando las keywords del usuario contra partidas del catálogo, la inferencia introdujo ambigüedad y ruido, degradando el ranking ($\Delta = -0{,}0135$). Por tanto, se descartó añadir complejidad algorítmica y columnas adicionales en base de datos.
\end{itemize}

\subsection{Decisiones Implementadas}

\begin{itemize}[leftmargin=*]
    \item \textbf{Implementación 1: Búsqueda Híbrida y Fusión por Rango Recíproco (\textit{Reciprocal Rank Fusion - RRF}).}
    Se implementó un canal léxico basado en BM25 sobre títulos, descripciones y partidas licitadas, fusionando sus clasificaciones con la búsqueda densa mediante la función agnóstica de escala:
    \begin{equation}
        RRF\_Score(d) = \sum_{m \in M} \frac{1}{k + \text{rank}_m(d) + 1}, \quad \text{con } k=60
    \end{equation}
    Esto eliminó la necesidad de calibrar pesos de combinación ad-hoc sensibles a la distribución de puntajes.

    \item \textbf{Implementación 2: Normalización Regional Canónica (\texttt{are\_regions\_matching}).}
    Se identificó y corrigió un error en el comparador de regiones: Mercado Público introduce espacios y variantes tipográficas (e.g., \texttt{'Región de Coquimbo '}), lo que hacía que comparaciones de igualdad estricta descartaran indebidamente $140$ de $250$ adjudicaciones válidas. Con normalización canónica, la región pasó de penalizar a ser un filtro y factor de afinidad de alta precisión.

    \item \textbf{Implementación 3: Fórmula Calibrada Multivariante ($R_c + B + C$).}
    Se reformuló el cálculo del porcentaje de compatibilidad como un modelo ordinal basado en dos regresiones logísticas:
    \begin{equation}
        P = 0{,}50 \cdot \sigma(z_{\text{relevante}}) + 0{,}50 \cdot \min\left(\sigma(z_{\text{exacto}}), \sigma(z_{\text{relevante}})\right)
    \end{equation}
    donde cada logit $z$ combina tres señales desacopladas de escala:
    \begin{equation}
        z = a + b \cdot \text{logit}(R_c) + c \cdot B + d \cdot C
    \end{equation}
    \begin{itemize}
        \item $R_c$: Score de cross-encoder alimentado con una consulta corta estructurada (``Rubro: ... Productos: ...'').
        \item $B$: \textit{Best match}, correspondiente al máximo coseno entre las palabras clave de la empresa y las partidas de la licitación.
        \item $C$: \textit{Coverage}, correspondiente al promedio sobre partidas del máximo coseno alcanzado.
    \end{itemize}

    \item \textbf{Implementación 4: Desacoplamiento del Reranker INT8.}
    Se eliminó el procesamiento en batch en el servicio \texttt{BgeRerankerService}, pasando a evaluar cada par individualmente. Esto corrigió el bug de cuantización dinámica y redujo el tiempo de procesamiento de $21$--$24$ segundos a tan solo $7$ segundos por cada 50 candidatas ($3\times$ de aceleración).

    \item \textbf{Implementación 5: Colección Léxica Dispersa en Qdrant (\texttt{tender\_lexical}) y Shadow Matching (ADR-0005).}
    Se implementó un repositorio desacoplado para vectores sparse indexados mediante tokenización propia (\texttt{LexicalTokenizer}) y cálculo de IDF en el servidor Qdrant (\texttt{Modifier.IDF}). Se integró una arquitectura de evaluación en sombra gobernada por \texttt{MATCHING\_SHADOW\_ENABLED} para auditar variaciones de ranking sin afectar a los usuarios en producción.
\end{itemize}

% =========================================================================
\section{Resultados Finales y Benchmarking}
% =========================================================================

\subsection{Comparativa de Estrategias sobre Métricas Duras}
La Tabla~\ref{tab:metricas_duras} resume los resultados obtenidos en el benchmark con split temporal sin fuga sobre los 50 proveedores y $1.176$ licitaciones del catálogo natural.

\begin{table}[htbp]
\centering
\small
\setlength{\tabcolsep}{4.5pt}
\caption{Comparativa General de Estrategias --- Métricas Duras ($N=50$, Catálogo = $1.176$).}
\label{tab:metricas_duras}
\begin{tabularx}{\textwidth}{lcccccccc}
\toprule
\textbf{Estrategia} & \textbf{NDCG@10} & \textbf{MRR} & \textbf{Succ@1} & \textbf{P@10} & \textbf{R@10} & \textbf{R@50} & \textbf{R@100} & \textbf{R@200} \\
\midrule
\textbf{A: Denso (BGE-M3)} & 0{,}0517 & 0{,}0675 & 0{,}0200 & 0{,}0160 & 0{,}0800 & 0{,}2400 & 0{,}3800 & 0{,}5300 \\
\textbf{B: Híbrido BM25 + RRF} & 0{,}0973 & 0{,}1225 & 0{,}0600 & 0{,}0300 & 0{,}1500 & 0{,}4500 & 0{,}5500 & 0{,}7500 \\
\textbf{C: Híbrido Estructural} & \textbf{0{,}1579} & \textbf{0{,}1816} & 0{,}0800 & \textbf{0{,}0480} & \textbf{0{,}2400} & \textbf{0{,}5900} & \textbf{0{,}6700} & \textbf{0{,}8000} \\
\textbf{E: Multi-Consulta} & 0{,}1193 & 0{,}1440 & 0{,}0600 & 0{,}0380 & 0{,}1900 & 0{,}4100 & 0{,}5500 & 0{,}7000 \\
\textbf{D: Reranker Top-100} & 0{,}1175 & 0{,}1632 & \textbf{0{,}1000} & 0{,}0320 & 0{,}1600 & 0{,}5200 & 0{,}6700 & 0{,}8000 \\
\bottomrule
\end{tabularx}
\end{table}

\begin{table}[htbp]
\centering
\small
\setlength{\tabcolsep}{5pt}
\caption{Pruebas de Significancia Estadística sobre NDCG@10 (Wilcoxon Pareado, Holm, Bootstrap 95\%).}
\label{tab:significancia}
\begin{tabularx}{\textwidth}{lccccc}
\toprule
\textbf{Comparación} & $\Delta$ \textbf{NDCG@10} & $\text{IC}_{95\%}$ & $p$-valor & $p$ \textbf{ajustado (Holm)} & \textbf{Conclusión} \\
\midrule
B vs A & $+0{,}0456$ & $[+0{,}0167, +0{,}0770]$ & 0{,}0087 & 0{,}0173 & Significativo ($\star$) \\
\textbf{C vs A} & $\mathbf{+0{,}1062}$ & $[\mathbf{+0{,}0572, +0{,}1600}]$ & \textbf{0{,}0004} & \textbf{0{,}0022} & \textbf{Altamente Significativo} ($\star\star$) \\
E vs A & $+0{,}0676$ & $[+0{,}0301, +0{,}1101]$ & 0{,}0028 & 0{,}0113 & Significativo ($\star$) \\
D vs A & $+0{,}0658$ & $[+0{,}0250, +0{,}1110]$ & 0{,}0031 & 0{,}0113 & Significativo ($\star$) \\
C vs B & $+0{,}0606$ & $[+0{,}0210, +0{,}1040]$ & 0{,}0089 & 0{,}0173 & Significativo ($\star$) \\
D vs C & $-0{,}0404$ & $[-0{,}0830, +0{,}0010]$ & 0{,}0630 & --- & No Significativo ($ns$) \\
\bottomrule
\end{tabularx}
\end{table}

\subsection{Evaluación Graduada con Juez Ciego (Pool Común de 1.146 Pares)}
Para verificar si las licitaciones no ganadas pero afines eran recomendadas adecuadamente, se conformó un pool de $1.146$ pares evaluados por modelos LLM independientes bajo acuerdo inter-juez sustancial ($\kappa$ ponderado $= 0{,}82$, acuerdo exacto del $80\%$).

\begin{table}[htbp]
\centering
\small
\caption{Evaluación con Juez LLM Ciego sobre Pool Común ($1.146$ pares evaluados).}
\label{tab:juez_llm}
\begin{tabularx}{\textwidth}{lcccc}
\toprule
\textbf{Estrategia} & \textbf{NDCG@10} & \textbf{P@10 ($rel \ge 1$)} & \textbf{P@10 ($rel = 2$)} & \textbf{Success@1 ($rel \ge 1$)} \\
\midrule
\textbf{A: Denso} & 0{,}3790 & 0{,}4900 & 0{,}2000 & 0{,}6200 \\
\textbf{B: Híbrido BM25} & 0{,}5790 & 0{,}7080 & 0{,}3140 & 0{,}9000 \\
\textbf{C: Híbrido Estructural} & \textbf{0{,}6600} & \textbf{0{,}7720} & \textbf{0{,}3800} & 0{,}9400 \\
\textbf{E: Multi-Consulta} & 0{,}6150 & 0{,}7260 & 0{,}3460 & 0{,}8800 \\
\textbf{D: Reranker Top-100} & 0{,}6000 & 0{,}7020 & 0{,}3180 & \textbf{0{,}9400} \\
\bottomrule
\end{tabularx}
\end{table}

\begin{quote}
\textit{Conclusión del Juez:} El orden relativo de las estrategias se preserva de manera idéntica ($C > E > D > B > A$), pero las brechas son notablemente mayores. La estrategia híbrida estructural alcanza un $77{,}2\%$ de precisión en el top-10 para licitaciones afines o exactas, frente a solo un $49{,}0\%$ del baseline denso.
\end{quote}

\subsection{Benchmarking del Pipeline Productivo y el Canal Léxico (Fase 0 y Fase 2)}
En la Fase 0 del plan técnico de despliegue, se comparó el pipeline real en producción ($P_{\text{Producción}}$, que implementa la recuperación bi-canal y la fórmula $R_c + B + C$) frente al agregado del canal léxico disperso ($P+L$).

\begin{table}[htbp]
\centering
\small
\caption{Benchmark del Pipeline de Producción Real vs. Canal Léxico ($P+L$).}
\label{tab:prod_lexico}
\begin{tabularx}{\textwidth}{lcccccc}
\toprule
\textbf{Variante} & \textbf{NDCG@10} & \textbf{MRR} & \textbf{Succ@1} & \textbf{P@10} & \textbf{Recall@50} & \textbf{Recall@100} \\
\midrule
$A_{\text{Denso}}$ (Referencia inicial) & 0{,}0503 & 0{,}0701 & 0{,}0200 & 0{,}0180 & 0{,}3500 & 0{,}4700 \\
$P_{\text{Producción}}$ (Producción actual) & 0{,}1409 & 0{,}2002 & \textbf{0{,}1400} & 0{,}0360 & 0{,}3500 & 0{,}4600 \\
$P+L_{\text{Scorer}}$ & 0{,}1400 & 0{,}1895 & 0{,}1200 & 0{,}0380 & 0{,}3500 & 0{,}6200 \\
$P+L_{\text{RRF}}$ (Producción + Sparse) & \textbf{0{,}1576} & \textbf{0{,}1995} & 0{,}1200 & \textbf{0{,}0440} & \textbf{0{,}4000} & \textbf{0{,}5000} \\
\bottomrule
\end{tabularx}
\end{table}

\begin{itemize}[leftmargin=*]
    \item El pipeline de producción actual ($P_{\text{Producción}}$) ya superaba con creces al baseline denso ($NDCG@10$: $0{,}1409$ vs. $0{,}0503$, $\Delta = +0{,}089$, $p=0{,}0015$).
    \item La incorporación del canal léxico fusionado mediante RRF ($P+L_{\text{RRF}}$) incrementa aún más la recuperación en rangos tempranos ($Recall@50$: $0{,}35 \to 0{,}40$, $Recall@100$: $0{,}46 \to 0{,}50$) y eleva el $NDCG@10$ a $0{,}1576$.
\end{itemize}

\subsection{Resultados de la Calibración de Compatibilidad}
La calibración mediante regresión logística multivariante ordinal ($R_c + B + C$) transformó radicalmente la confiabilidad del score final presentado al usuario:

\begin{table}[htbp]
\centering
\small
\caption{Métricas de Calibración Probabilística y Capacidad de Separación ($1.146$ pares).}
\label{tab:calibracion}
\resizebox{\textwidth}{!}{
\begin{tabular}{lccccc}
\toprule
\textbf{Modelo} & \textbf{AUC-ROC} & \textbf{Brier Score} & \textbf{ECE} & \textbf{Precisión en $\ge 70\%$} & \textbf{Media Ordinal (Ajena / Afín / Exacta)} \\
\midrule
Fórmula Original ($P_0$) & 0{,}5827 & 0{,}3266 & 0{,}2828 & --- & $27\% \;/\; 30\% \;/\; 33\%$ (Plana/Comprimida) \\
Reranker solo ($R$) & 0{,}5402 & 0{,}2444 & 0{,}0213 & $0{,}0\%$ & $40\% \;/\; 41\% \;/\; 41\%$ \\
Reranker corto ($R_c$) & 0{,}6402 & 0{,}2302 & 0{,}0332 & $100\%$ ($1{,}2\%$ volumen) & $37\% \;/\; 41\% \;/\; 46\%$ \\
Best Match + Coverage ($B+C$) & 0{,}7615 & 0{,}2043 & 0{,}0775 & $92{,}9\%$ & $35\% \;/\; 42\% \;/\; 49\%$ \\
\textbf{Calibrado ($R_c + B + C$)} & \textbf{0{,}7489} & \textbf{0{,}2024} & \textbf{0{,}0575} & \textbf{97{,}6\%} & $\mathbf{33{,}5\% \;/\; 41{,}9\% \;/\; 51{,}7\%}$ \\
\bottomrule
\end{tabular}
}
\end{table}

\begin{itemize}[leftmargin=*]
    \item \textbf{AUC-ROC saltó de $0{,}58$ a $0{,}75$:} El sistema adquiere capacidad efectiva de discernimiento entre licitaciones ajenas y relevantes.
    \item \textbf{Brier Score y ECE reducidos a mínimos históricos:} El error de calibración esperado (\textit{Expected Calibration Error, ECE}) se redujo de $0{,}2828$ a $0{,}0575$, garantizando que una predicción del $70\%$ corresponda empíricamente a un $\sim 70\%$ de aciertos.
    \item \textbf{Selector Verde Operativo:} Cuando la fórmula calibrada arroja un puntaje $\ge 70\%$, la precisión real verificada de que la licitación es pertinente es del \textbf{97{,}6\%} (frente a $0\%$ en la fórmula original, que jamás alcanzaba ese rango).
\end{itemize}

% =========================================================================
\section{Conclusiones y Recomendaciones de Ingeniería}
% =========================================================================

\begin{enumerate}[leftmargin=*]
    \item \textbf{Superación del paradigma denso puro:} En dominios técnicos y compras públicas, los vectores densos generales son insuficientes por sí solos ($Recall@100 = 38\%$). La combinación sinérgica de búsqueda semántica densa, índices léxicos dispersos (BM25) y filtros estructurales normalizados es el único diseño validado empíricamente capaz de triplicar el $NDCG@10$ con significancia estadística.
    \item \textbf{El primer estadio define el límite del sistema:} El uso de rerankers pesados en estadios iniciales no compensa un primer estadio con bajo Recall. La inversión computacional debe priorizar la exhaustividad de candidatos antes de añadir capas de reordenamiento profundo.
    \item \textbf{La calibración probabilística es mandataria:} Ningún score heurístico o combinación lineal intuitiva garantiza consistencia operativa. La adopción de la fórmula calibrada $R_c + B + C$ permitió dignificar la experiencia de usuario, habilitando con rigor los umbrales de alerta y selector verde en Chiripa.
    \item \textbf{Validación continua en sombra:} La arquitectura de \textit{Shadow Matching} establecida mediante el ADR-0005 y la recolección de eventos de telemetría de usuario garantizan que cualquier iteración algorítmica futura se evalúe contra clics y postulaciones reales antes de impactar el tráfico productivo.
\end{enumerate}

\end{document}
